% 第 5 题程序框图：按固定源内容流矢量层精确值复刻（1 pt → 0.03515 cm，与全库同制）。
% 原点取框图竖轴 x=233.4258 pt、底缘 y（顶向下）=764.9258 pt；y 向上。
% 源实测：7 框（开始/结束=圆角 27.988 宽、框2=矩形 96.348×28.063 含真分式、框3 37.148×13.555、
% 框4 50.656×15.160、菱形 126.215×31.430、输出=平行四边形（斜偏 6.1328））；
% 6 段竖直连接 + 「否」支 3 段直角回边（左行 19.836、上行 80.930、右行 79.641，箭头落 框2/框3 之间竖线上）。
\begin{tikzpicture}[x=0.03515cm,y=0.03515cm,>=stealth,every node/.style={font=\small,inner sep=0pt},
  gl/.style={line width=0.398bp,line cap=round,line join=round}]
% ===== 7 框 =====
\draw[gl,rounded corners=0.1401cm] (-13.9961,209.6367) rectangle (13.9922,226.4688);   % 1 开始
\draw[gl] (-48.1758,167.0078) rectangle (48.1718,195.0703);                            % 2 k=0,a=3,q=1/2
\draw[gl] (-18.5742,138.8868) rectangle (18.5742,152.4414);                            % 3 a=aq
\draw[gl] (-25.3281,109.1602) rectangle (25.3281,124.3203);                            % 4 k=k+1
\draw[gl] (63.1055,78.7969)--(0,94.5117)--(-63.1094,78.7969)--(0,63.0821)--cycle;      % 5 菱形 a<1/4
\draw[gl] (-24.7774,31.5781)--(-18.6446,48.4297)--(24.7695,48.4297)--(18.6367,31.5781)--cycle; % 6 输出 k
\draw[gl,rounded corners=0.1401cm] (-13.9961,0) rectangle (13.9922,17.0117);           % 7 结束
% ===== 竖直连接（6 段，均 x=0）=====
\draw[gl,->] (0,209.4375)--(0,195.2696);   % 开始→框2
\draw[gl,->] (0,166.8125)--(0,152.6446);   % 框2→框3
\draw[gl,->] (0,138.6875)--(0,124.5196);   % 框3→框4
\draw[gl,->] (0,108.9610)--(0,94.7930);    % 框4→菱形顶
\draw[gl,->] (0,62.8008)--(0,48.6328);     % 菱形底→框6
\draw[gl,->] (0,31.3789)--(0,17.2110);     % 框6→结束
% ===== 「否」回边（3 段直角折线 + 箭头，落 框2/框3 之间竖线）=====
\draw[gl,->] (-63.3906,78.7969)--(-83.2266,78.7969)--(-83.2266,159.7266)--(0,159.7266);
% ===== 框内文字与分支标签（照源字形包围盒中心定位）=====
\node at (0,218.5358) {开始};
\node at (0,181.0390) {\(k=0,\ a=3,\ q=\dfrac{1}{2}\)};
\node at (0,146.9200) {\(a=aq\)};
\node at (0,116.0158) {\(k=k+1\)};
\node at (0,78.0860) {\(a<\dfrac{1}{4}\)};
\node at (-0.0800,40.1658) {输出 \(k\)};
\node at (0,8.9358) {结束};
\node at (-63.3258,87.6058) {否};
\node at (9.0342,56.2358) {是};
\end{tikzpicture}
